\chapter{Conventions}
\label{ch:conventions}

Most silent errors in an AT-TPC analysis come from mixing coordinate frames, field signs or units. This chapter
collects the conventions used in the framework. Where a convention was measured on a particular data set rather
than fixed by the code, it says so; such values must be re-measured for a new set-up.

\section{Coordinate frames}

\paragraph{Pad plane and beam axis.} $x$ and $y$ lie in the pad plane; the beam travels along $z$, on or near the
axis $x = y = 0$. Pad positions come from the pad map (\code{AtMap::CalcPadCenter}).

\paragraph{Digitisation frame ($z_\mathrm{digi}$).} The pulse-shape analysis converts the time bucket of a pulse into
the \emph{drift length}: \code{AtPSA::CalculateZGeo} returns
\begin{equation}
  z_\mathrm{digi} = Z_\mathrm{pad\,plane} - (\mathrm{TB}_\mathrm{entrance} - t_\mathrm{peak})\,
  \Delta t_\mathrm{TB}\, v_\mathrm{drift},
\end{equation}
with \code{ZPadPlane}, \code{TBEntrance} and the drift velocity from \code{AtDigiPar}. $z_\mathrm{digi} = 0$ at the pad
plane and $z_\mathrm{digi} = Z_\mathrm{pad\,plane}$ at the far end of the drift volume. Hits (\code{AtHit}) and track
points are stored in this frame.

\paragraph{Lab frame ($z_\mathrm{lab}$).} Physics quantities -- vertex position, polar angle -- are expressed along
the beam: $z_\mathrm{lab} = Z_\mathrm{pad\,plane} - z_\mathrm{digi}$. The fitters apply this transform
(\code{SetZPadPlane} on \code{AtGenfitter}; \code{SetZTransform(-1, ZPadPlane)} on \code{AtHelix3D}, through
\code{AtPIDTask::SetZPadPlane}). This value is set separately from the parameter file and the two need not agree:
\code{ATTPC.a1975\_deuterium\_dv1104.par} has \code{ZPadPlane} = 971.7312~mm, which follows from its drift velocity
and two time-bucket anchors, while \code{fitGenfitter\_a1975\_deuterium.C} calls \code{SetZPadPlane(1000.0)} and the
default transform of \code{AtHelix3D} also uses 1000. A difference is a constant offset of $z_\mathrm{lab}$: it does not
change angles or momenta, but it does shift the vertex position. Check which value applies when comparing vertices
between chains or with a beam energy-loss calculation. Internally, AtHelix3D's centroids are in the transformed frame.

\paragraph{Spyral's $z$ branch.} \code{AtPSA} can instead use a two-point $z$ calibration copied from Spyral
(\code{SetSpyralZ}). That branch returns 0 at the beam \emph{window} and the detector length at the Micromegas -- the
opposite of \code{CalculateZGeo}. A chain that enables it and then applies $z_\mathrm{lab} = Z_\mathrm{pad\,plane} -
z_\mathrm{digi}$ gets a mirrored $z_\mathrm{lab}$. No production macro enables it.

\paragraph{Reversed detector.} When the beam enters through the pad plane (\code{ReverseDrift}), the digitisation
keeps $z_\mathrm{digi} = Z_\mathrm{pad\,plane} - z_\mathrm{beam}$, so the hit cloud and every downstream convention
are unchanged; only the drift length of a given point differs. See Chapter~\ref{ch:parameters}.

\section{Directions and angles}

\paragraph{Polar angle.} $\theta$ is measured from the beam direction ($+z_\mathrm{lab}$). A track is
\emph{backward} when its vertex end has the higher $z_\mathrm{lab}$.

\paragraph{Spyral direction.} \code{AtSpyralResult::direction} is 0 (forward) when the vertex is at the
\emph{lowest} $z_\mathrm{digi}$, 1 (backward) otherwise, $-1$ when undecided.

\paragraph{Pattern-recognition direction.} The forward/backward call of \code{AtPRA} (\code{GeoTheta}) comes from the
sign of one fitted line; see Chapter~\ref{ch:patternrecognition}. AtHelix3D decides the vertex end by a two-of-three
vote instead (Chapter~\ref{ch:pid}).

\paragraph{Rotation sense.} AtHelix3D's rotation-sense vote expects the track, travelling inward from the vertex end,
to turn with a given sense in the digitisation $x$--$y$ plane (\code{SetExpectedSense}, default $+1$). On a1975 this
was measured as $+1$ for both data and simulation at 2.85~T. It must be re-measured for another field, charge sign or
geometry.

\section{Magnetic-field sign}

\code{AtDigiPar} carries only the magnitude of the field. The sign is set where it is needed:
\begin{itemize}
\item in the simulation, through the field object of the transport macro (for example
  \code{AtConstField::SetField(0, 0, bFieldkG)} with $-28.5$~kG in \code{C16\_dt\_sim.C}, generated in the data
  convention so that tracks turn the same way as in the data);
\item in the fitters, through the field argument (\code{AtGenfitter}, $-2.85$ for the a1975 data) or
  \code{SetBFieldSign} (\code{AtMCQMinimization} and the UKF chain).
\end{itemize}
The wrong sign does not always fail loudly: a fit can converge to a wrong momentum, and in \code{AtMCQMinimization} a
simulated track curls away from the data and no set of parameters brings it back. Check the sign first when fits look
wrong. The cases measured so far are tabulated in Section~\ref{sec:fit-frames}.

On the a1975 simulation, the reconstructed polar angle of Spyral came out as $180^\circ - \theta_\mathrm{true}$ and
the field sign of the fitters had to be $+1$ for that simulation and $-1$ for the data; the sign belongs to the pair of
simulation field and fitter, not to ``simulation'' as such. Matching reconstructed angles
directly against truth without accounting for this gives zero acceptance, which is how the effect shows up. Measure
the relation between reconstructed and true angles on the simulation of each new set-up before using it.

\section{Units}

\begin{center}
\small
\begin{tabular}{p{0.32\textwidth}p{0.6\textwidth}}
\toprule
Quantity & Unit \\
\midrule
lengths, positions & mm (hits, tracks, fitters); cm in \code{AtDigiPar} drift velocity (cm/$\mu$s) and in Geant4
  geometry \\
energies & MeV (MeV per nucleon where a generator says so) \\
magnetic field & tesla in the reconstruction; kG in \code{AtConstField} \\
rigidity & T\,m \\
gas density & kg/m$^3$ in \code{AtDigiPar} (numerically equal to mg/cm$^3$); g/cm$^3$ in \code{AtELossCATIMA} --
  a factor 1000 between the two \\
time & time buckets in traces; $\mu$s for peaking time and \code{TBTime} \\
\bottomrule
\end{tabular}
\end{center}

\section{Masses}

\paragraph{Generator masses.} The AT-TPC simulation macros pass ion masses to \code{AtTPCIonGenerator} in atomic
mass units although the underlying \code{FairIon} parameter is documented in GeV. The convention is consistent across
the macros; verify the beam energy from the Monte Carlo truth when setting up a new simulation.

\paragraph{Atomic and nuclear masses.} In a two-body Q value the electron masses cancel only if the convention is
consistent: with \emph{atomic} masses for the heavy particles, the light particles must also be atomic, or all four
must be nuclear. Mixing them shifts the Q value by 0.511~MeV per unbalanced electron, which reads as a calibration
offset. The test costs nothing: compute Q for elastic scattering from the masses in the code; it must be exactly zero.

\section{Track identifiers and order}

\begin{itemize}
\item Match PID results, fits and pattern tracks on \code{trackID}, never on array positions.
\item Hits in a pattern track come in drift-time order. Within one time bucket the order is arbitrary, which
  matters for nearly flat tracks and for fitters that rely on the hit order
  (Chapters~\ref{ch:patternrecognition} and \ref{ch:fitting}).
\end{itemize}
